\documentclass[12pt]{article}

% Match weeks 2--5 using fonts available in the local TeX installation.
\usepackage[T1]{fontenc}
\usepackage{times}
\usepackage{microtype}
\usepackage{amsmath}
\usepackage{amsfonts} % for \mathbb{E}
\usepackage[parfill]{parskip}
\setlength{\parindent}{0pt}
\setlength{\parskip}{2ex plus 0.2ex minus 0.2ex}
\usepackage[letterpaper,margin=1in]{geometry}
\usepackage[hidelinks]{hyperref}

\begin{document}

\begin{flushleft}
\textbf{Week 6: Introduction to Variational Inference} \\
Jaisel Singh (js6897) \\
October 21, 2026
\end{flushleft}

\vspace{0.1in}

What stood out to me in Blei, Kucukelbir, and McAuliffe's review was that
inference becomes an optimization problem rather than a sampling problem. We
choose a family of densities over the latent variables and search within it for the
member closest to the posterior in KL divergence. The obstacle is that
this divergence contains the log evidence, the quantity we could not compute in
the first place, so we maximize the evidence lower bound instead. Because the
ELBO equals the log evidence minus that divergence, maximizing it minimizes the
divergence without ever evaluating the constant. I found the second form,
\mbox{$\mathbb{E}[\log p(x \mid z)] - \mathrm{KL}(q(z) \| p(z))$}, the clearest
way to see the tradeoff: fit to the observed data against staying close to the
prior. It also makes week 3's stochastic optimization directly relevant, since the
ELBO is one more objective to climb.

The cost of this convenience is where the review was most useful to me. The
mean-field family treats the latent variables as mutually independent, each with
its own factor, which makes coordinate ascent cheap but removes any ability to
represent dependence among them. The authors are direct that variational
inference generally underestimates posterior variance, a consequence of the
objective rather than a fixable implementation detail. They also say
plainly that the relative accuracy of variational inference and MCMC is still not
understood. That contrast sharpens week 4 for me: MCMC pays in computation for
samples that are asymptotically exact, while this pays in accuracy for speed, and
the value of the ELBO at convergence does not report how large that error is.

Jordan, Ghahramani, Jaakkola, and Saul arrive at the same methods from an older
direction. Their emphasis is convex duality: a concave function can be written as
a minimum over a family of linear upper bounds, as with
\mbox{$\ln x \le \lambda x - \ln \lambda - 1$}, where each choice of the
variational parameter $\lambda$ is tight at one point. Fitting those parameters
then turns an intractable graphical model into a simpler one where inference runs
efficiently, and what comes back is a bound on the quantity of interest. Reading
the two together, the question I am left with is how to check whether the
independence assumption has distorted the particular posterior summary I care
about, given that the bound itself will not tell me.

\vspace{0.05in}
{\footnotesize
\textit{Readings:} Blei et al.\ (2017), \href{https://arxiv.org/abs/1601.00670}{\emph{Variational Inference: A Review for Statisticians}};
Jordan et al.\ (1999), \href{https://www.cs.columbia.edu/~blei/fogm/readings/JordanGhahramaniJaakkolaSaul1999.pdf}{\emph{An Introduction to Variational Methods for Graphical Models}}.
\par}

\end{document}

%%% Local Variables:
%%% mode: latex
%%% TeX-master: t
%%% End:
